% Part: first-order-logic
% Chapter: axiomatic-deduction
% Section: proof-theoretic-notions

\documentclass[../../../include/open-logic-section]{subfiles}

\begin{document}

\iftag{FOL}
      {\olfileid{fol}{axd}{ptn}}
      {\olfileid{pl}{axd}{ptn}}

\olsection{Begrippen over formele afleidingen}

\begin{explain}
We hebben al een paar belangrijke begrippen over betekenis
gedefinieerd: \iftag{FOL}{geldigheid}{tautologie}, logisch gevolg en
vervulbaarheid, dus de vraag of de betrokken formules samen waar te
maken zijn. Nu definiëren we de bijbehorende
\emph{bewijstheoretische begrippen}: begrippen over formele
afleidingen. Daarbij vragen we niet of !!{sentence}s waar zijn in
!!{structure}s. We vragen of bepaalde formules wel of niet
!!{derivable} zijn. Een belangrijke ontdekking was dat de twee soorten
begrippen samenvallen. De \emph{correctheids-} en
\emph{volledigheidsstellingen} zeggen precies dat dit zo is.
\end{explain}

\begin{defn}[!!^{derivability}]
We zeggen dat !!a{formula}~$!A$ \emph{!!{derivable} is uit} $\Gamma$.
Dit schrijven we als $\Gamma \Proves !A$. Het betekent dat er
!!a{derivation} uit~$\Gamma$ bestaat waarvan~$!A$ de laatste formule
is.
\end{defn}

\begin{defn}[Stellingen]
!!^a{formula}~$!A$ is een \emph{stelling} als er !!a{derivation} van
$!A$ uit de lege verzameling bestaat, dus zonder gegeven formules als
aannamen. We schrijven $\Proves !A$ als $!A$ een stelling is. Als dat
niet zo is, schrijven we $\Proves/ !A$.
\end{defn}

\begin{defn}[Consistentie]
Een verzameling !!{formula}s~$\Gamma$ is precies dan
\emph{consistent} als $\Gamma\Proves/ \lfalse$. Er is dan geen
afleiding van een tegenspraak uit die verzameling. Anders is de
verzameling \emph{inconsistent}.
\end{defn}

\begin{prop}[Reflexiviteit: wat is aangenomen, volgt ook]
\ollabel{prop:reflexivity}
Als $!A \in \Gamma$, dan $\Gamma \Proves !A$.
\end{prop}

\begin{proof}
  De rij die alleen uit de !!{formula}~$!A$ bestaat, is al
  !!a{derivation} van~$!A$ uit~$\Gamma$.
\end{proof}

\begin{prop}[Monotoniciteit: extra aannamen mogen erbij]
\ollabel{prop:monotonicity}
Als $\Gamma \subseteq \Delta$ en $\Gamma \Proves !A$, dan $\Delta
\Proves !A$.
\end{prop}

\begin{proof}
Elke !!{derivation} van $!A$ uit $\Gamma$ is ook !!a{derivation} van
$!A$ uit~$\Delta$.
\end{proof}

\begin{prop}[Transitiviteit: afleidingen achter elkaar zetten]
\ollabel{prop:transitivity}
Als $\Gamma \Proves !A$ en $\{!A\} \cup \Delta \Proves
!B$, dan $\Gamma \cup \Delta \Proves !B$.
\end{prop}

\begin{proof}
  Stel dat $\{!A\} \cup \Delta \Proves !B$. Dan is er
  !!a{derivation} $!B_1$, \dots, $!B_l = !B$ uit~$\{!A\} \cup
  \Delta$. Sommige stappen in die !!{derivation} gebruiken via een
  regel een eerdere regel~$!B_i = !A$. We weten ook dat er
  !!a{derivation} van~$!A$ uit~$\Gamma$ bestaat. Dat is een
  !!a{derivation}~$!A_1$, \dots, $!A_k = !A$, waarin elke $!A_i$ een
  axioma, !!a{element} van~$\Gamma$ of een juiste toepassing van een
  afleidingsregel is. Zet de twee rijen nu achter elkaar:
  \[
  !A_1, \dots, !A_k = !A, !B_1, \dots, !B_l = !B.
  \]
  Dit is een juiste !!{derivation} van~$!B$ uit $\Gamma \cup
  \Delta$. Elke regel met $B_i = !A$ krijgt nu dezelfde
  rechtvaardiging als~$!A_k = !A$.
\end{proof}

Hieruit volgen twee nuttige gevallen. Als $\Gamma \Proves !A$ en $!A
\Proves !B$, dan $\Gamma \Proves !B$. En als $!A_1,
\dots, !A_n \Proves !B$ en $\Gamma \Proves !A_i$ voor elke~$i$, dan
$\Gamma \Proves !B$.

\begin{prop}
\ollabel{prop:incons}
$\Gamma$ is precies dan inconsistent als $\Gamma \Proves !A$ geldt
voor elke~$!A$.
\end{prop}

\begin{proof}
Oefening.
\end{proof}

\tagprob{FOL}
\begin{prob}
Bewijs \olref[fol][axd][ptn]{prop:incons}.
\end{prob}
\tagendprob

\tagprob{notFOL}
\begin{prob}
Bewijs \olref[pl][axd][ptn]{prop:incons}.
\end{prob}
\tagendprob

\begin{prop}[Compactheid: eindige delen zijn genoeg]
\ollabel{prop:proves-compact}
  \begin{enumerate}
  \item Als $\Gamma \Proves !A$, dan is er een eindige
    deelverzameling $\Gamma_0 \subseteq \Gamma$ waarvoor
    $\Gamma_0 \Proves !A$.
  \item Als elke eindige deelverzameling van~$\Gamma$ consistent is,
    dan is $\Gamma$ zelf consistent.
  \end{enumerate}
\end{prop}

\begin{proof}
  \begin{enumerate}
    \item Als $\Gamma \Proves !A$, is er een eindige rij
      !!{formula}s $!A_1$, \dots,~$!A_n$ met $!A \ident !A_n$. Elke
      $!A_i$ is een logisch axioma, !!a{element} van~$\Gamma$, of volgt
      met modus ponens uit eerdere !!{formula}s. Laat $\Gamma_0$
      bestaan uit de $!A_i$ die in~$\Gamma$ staan. Dezelfde
      !!{derivation} is dan ook !!a{derivation} uit~$\Gamma_0$. Dus
      $\Gamma_0 \Proves !A$.
    \item Dit is de contrapositie van~(1): dezelfde bewering met de
      richting omgekeerd en beide kanten ontkend. Hier gebruiken we
      het bijzondere geval $!A \ident \lfalse$.
\end{enumerate}
\end{proof}

\end{document}
